\documentclass[10pt,a4paper]{amsart}
\usepackage[margin=19mm]{geometry}
\usepackage{amsmath,amssymb,mathtools}
\usepackage[scr=rsfs,cal=euler]{mathalfa}
\usepackage{xfrac}
\usepackage[unicode,hidelinks,bookmarks=false]{hyperref}
\newcommand{\ie}{i.e.\ }
\newcommand{\cf}{cf.\ }
\newcommand{\resp}{respectively\ }
\newcommand{\Subsection}{\subsection}
\providecommand{\nobreakdash}{\nobreak\hbox{-}}
\setlength{\emergencystretch}{2em}
\multlinegap0pt
\usepackage{bm}
\usepackage{amssymb}
\usepackage{tikz}
\usepackage{xcolor}
\usepackage{algorithm}\usepackage{algpseudocode}
\usepackage{float}
\newtheorem{theorem}{Theorem}
\title{Wasserstein projection estimators for circular distributions}
\author{Naoki Otani, Takeru Matsuda}
\date{Source: arXiv:2510.18367v2 (CC BY 4.0)}
\begin{document}
\maketitle

\noindent\emph{Source and licence.} Wasserstein projection estimators for circular distributions. Version arXiv:2510.18367v2 (CC BY 4.0), distributed by arXiv under the Creative Commons Attribution 4.0 International licence, http://creativecommons.org/licenses/by/4.0/, as recorded in the arXiv licence metadata for this exact version and shown on the abstract page https://arxiv.org/abs/2510.18367v2. Full licence notice: \texttt{licenses/arxiv-2510.18367-CC-BY-4.0.txt}.\par\smallskip

\noindent\textit{Editorial scope.} Counted unit: the article's theorem theorem (source0.tex lines 210--212). The article's complete original proof of it is delivered with it (source0.tex lines 213--234, one \texttt{proof} environment of 22 lines). No mathematics is written, repaired or re-derived: the statement and the proof are the article's own lines, in the article's own notation, and no retained line is substituted or re-flowed.  No further source-local context is needed: every cross-reference of every retained line resolves inside the two delivered spans.  Nothing was omitted from any retained span, so the package carries no omission record.  The retained lines cite 2 works, whose entries are transcribed verbatim from the article's own bibliography source in the e-print and delivered with short editorial labels recorded in the item's reference mapping.  No retained line contains a figure environment, an image-inclusion command or drawing code, so no image is delivered and the package declares no asset dependency.  Editorial formatting only, in addition: an AMS \texttt{amsart} preamble that restates the article's own declarations and macros used by the retained lines, namely the environments \texttt{theorem} on separate counters, together with the standard packages those lines need.  The retained environments print in this document as Theorem 1 in order of appearance, whereas the article prints the same items with the numbering of its own full document.  The complete native source of the article as distributed in the arXiv e-print for this exact version is bundled unchanged as \texttt{sources/187496/source0.tex}.
\begin{theorem}
	If $p \geq 1$ and the model is parametric and identifiable, then $\hat{\theta}_{W,p} \to \theta$ as $n \to \infty$.
\end{theorem}

\begin{proof}
	From the triangle inequality,
	\[
	|W_p(\hat{q},p_{\theta}) - W_p(p_{\theta_0},p_{\theta})| \leq W_p(\hat{q},p_{\theta_0})
	\]
	for every $\theta$.
	From \cite{Fournier},
	\[
	W_p(\hat{q},p_{\theta_0}) \to 0
	\]
	as $n \to \infty$.
	Thus,
	\[
	\sup_{\theta} |W_p(\hat{q},p_{\theta}) - W_p(p_{\theta_0},p_{\theta})| \to 0
	\]
	as $n \to \infty$.
	Also, since the model is parametric and  identifiable, for every $\varepsilon>0$,
	\[
	\inf_{\theta: \| \theta-\theta_0 \| > \varepsilon} W_p(p_{\theta_0},p_{\theta}) > 0 = W_p(p_{\theta_0},p_{\theta_0}),
	\]
	Therefore, from Theorem~5.7 of \cite{vV}, $\hat{\theta} \to \theta_0$.
\end{proof}

\begin{thebibliography}{99}
\bibitem[Fourni]{Fournier} Fournier, N. \& Guillin, A. (2015). On the rate of convergence in Wasserstein distance of the empirical measure. \textit{Probability Theory and Related Fields}, \textbf{162}, 707--738.
\bibitem[vV]{vV} van der Vaart, A. W. (1998). \textit{Asymptotic Statistics}. Cambridge University Press.
\end{thebibliography}
\end{document}
